% Einheit 2 von 4 – Zeit (bank/einheiten/e2.jsonl, zusammenbau v0.1)
%% TODO Zweigzeile Teil 1: was hier gelernt wird (Satz in Schülersprache aus dem Einheitstitel) – Einheit 2
\einheitenkopf[e2]{Einheit 2 von 4 · Zeit}
\zweigzeile{ab Kl. 5 · P10}
\setcounter{aufgabe}{22}

%% TODO Titel: Ich-kann-Satz fehlt in der Bank (Platzhalter: Kettenname) – Nr. 23
%% TODO Anweisung: ein Satz über der ersten Teilaufgabe fehlt in der Bank – Nr. 23
\begin{aufgabe}{Zeit umrechnen und Zeitspannen berechnen}
\begin{teile}
\teil Der Film beginnt um 19:15 Uhr. Zeitpunkt oder Zeitspanne? Kreuze an.\\ \kreuz{Zeitpunkt}\\ \kreuz{Zeitspanne}
\end{teile}
%% TODO Darstellung für schwach fehlt in der Bank (einheiten-e2-k1-s1-v1; Abschnitt „Für schwache Schüler“)
\swz{$3$ min – wie viele s?}{}
%% TODO Darstellung für schwach fehlt in der Bank (einheiten-e2-k1-s1-v2; Abschnitt „Für schwache Schüler“)
\swz{$2$ h – wie viele min?}{}
%% TODO Darstellung für schwach fehlt in der Bank (einheiten-e2-k1-s1-v3; Abschnitt „Für schwache Schüler“)
\swz{$5$ min – wie viele s?}{}
%% TODO Darstellung für schwach fehlt in der Bank (einheiten-e2-k1-s1-v4; Abschnitt „Für schwache Schüler“)
\swz{$4$ h – wie viele min?}{}
%% TODO Darstellung für schwach fehlt in der Bank (einheiten-e2-k1-s1-v5; Abschnitt „Für schwache Schüler“)
\swz{$7$ min – wie viele s?}{}
\swfrage{Was bleibt gleich, was ändert sich?}
%% TODO Darstellung für schwach fehlt in der Bank (einheiten-e2-k1-s2-v1; Abschnitt „Für schwache Schüler“)
\swz{$240$ s – wie viele min?}{}
%% TODO Darstellung für schwach fehlt in der Bank (einheiten-e2-k1-s3-v1; Abschnitt „Für schwache Schüler“)
\swz{Eine Viertelstunde – wie viele min?}{}
%% TODO Darstellung für schwach fehlt in der Bank (einheiten-e2-k1-s4-v1; Abschnitt „Für schwache Schüler“)
\swz{$1{,}25$ h – wie viele min?}{}
%% TODO Darstellung für schwach fehlt in der Bank (einheiten-e2-k1-s5-v1; Abschnitt „Für schwache Schüler“)
\swz{$90$ min – wie viele h?}{}
%% TODO Darstellung für schwach fehlt in der Bank (einheiten-e2-k1-s6-v1; Abschnitt „Für schwache Schüler“)
\swz{Von 8:15 Uhr bis 9:40 Uhr – wie lange?}{}
%% TODO Darstellung für schwach fehlt in der Bank (einheiten-e2-k1-s7-v1; Abschnitt „Für schwache Schüler“)
\swz{Von 9:45 Uhr bis 12:15 Uhr – wie lange?}{}
%% TODO Darstellung für schwach fehlt in der Bank (einheiten-e2-k1-s8-v1; Abschnitt „Für schwache Schüler“)
\swz{Tim fährt um 8:30 Uhr los, fährt $1$ h $45$ min, macht $25$ min Pause und fährt dann noch $55$ min. Wann kommt er an?}{}
%% TODO Darstellung für schwach fehlt in der Bank (einheiten-e2-k1-s9-v1; Abschnitt „Für schwache Schüler“)
\swz{Eine Radtour dauert $5{,}5$ h. Wie viele Minuten sind das? \hfill (P10 2018 OS)}{}
\swz[3]{Ein Zug fährt um 9:47 Uhr ab und ist $3$ h $25$ min unterwegs. Wann kommt er an? \hfill (P10 2021 OS)}{}
\end{aufgabe}

%% TODO Titel: Ich-kann-Satz fehlt in der Bank (Platzhalter: Kettenname) – Nr. 24
%% TODO Anweisung: ein Satz über der ersten Teilaufgabe fehlt in der Bank – Nr. 24
\begin{aufgabe}{Zeiteinheiten nennen und ordnen}
%% TODO Darstellung für schwach fehlt in der Bank (einheiten-e2-k2-s1-v1; Abschnitt „Für schwache Schüler“)
\swz{Ordne, mit der kleinsten beginnend: Tag, Sekunde, Stunde, Woche, Minute}{}
\end{aufgabe}

%% TODO Titel: Ich-kann-Satz fehlt in der Bank (Platzhalter: Kettenname) – Nr. 25
%% TODO Anweisung: ein Satz über der ersten Teilaufgabe fehlt in der Bank – Nr. 25
\begin{aufgabe}{Startzeit aus Endzeit und Dauer}
%% TODO Darstellung für schwach fehlt in der Bank (einheiten-e2-k3-s1-v1; Abschnitt „Für schwache Schüler“)
\swz{Der Unterricht endet um 15:25 Uhr und dauert mit Pausen $6$ h $40$ min. Wann beginnt er?}{}
\end{aufgabe}

%% TODO Titel: Ich-kann-Satz fehlt in der Bank (Platzhalter: Kettenname) – Nr. 26
%% TODO Anweisung: ein Satz über der ersten Teilaufgabe fehlt in der Bank – Nr. 26
\begin{aufgabe}{Fahrplan: nächste Abfahrt, Fahrzeit}
\swa{Mia ist um 7:25 Uhr an der Haltestelle Rathaus. Mit welchem Bus fährt sie, und wie lange braucht sie bis zur Schule? Bus \leerfeld , \leerfeld[min]}{\sachtabelle{lccc}{Halt & Bus 1 & Bus 2 & Bus 3}{Rathaus & 7:12 & 7:42 & 8:12\\ Bahnhof & 7:31 & 8:01 & 8:31\\ Schule & 7:48 & 8:18 & 8:48}}
\end{aufgabe}

%% TODO Titel: Ich-kann-Satz fehlt in der Bank (Platzhalter: Kettenname) – Nr. 27
%% TODO Anweisung: ein Satz über der ersten Teilaufgabe fehlt in der Bank – Nr. 27
\begin{aufgabe}{Dauer in eine sinnvolle Einheit bringen (Sekunden in Minuten, Stunden in Tage)}
%% TODO Darstellung für schwach fehlt in der Bank (einheiten-e2-k5-s1-v1; Abschnitt „Für schwache Schüler“)
\swz{$72$ h – wie viele Tage?}{}
\end{aufgabe}

%% TODO Titel: Ich-kann-Satz fehlt in der Bank (Platzhalter: Kettenname) – Nr. 28
%% TODO Anweisung: ein Satz über der ersten Teilaufgabe fehlt in der Bank – Nr. 28
\begin{aufgabe}{Zeit umrechnen und Zeitspannen berechnen – Fehler finden, Begründen, Anwendung, Darstellungswechsel}
\swz[3]{Ben schreibt: $2{,}75$ h $= 2$ h $75$ min. Wo ist der Fehler?}{}
\swfrage{Begründe: Warum sind $0{,}4$ h nicht $40$ min?}
\swz[3]{Ein Handy-Akku lädt in $1{,}75$ h voll. Jonas steckt das Handy um 17:40 Uhr an und muss um 19:15 Uhr los. Ist der Akku dann voll?}{}
\swa{Ergänze die Tabelle für $1{,}75$ h.}{\sachtabelle{ccc}{Dezimalstunde & Stunden und Minuten & Minuten}{1,75 h & \leerzelle & \leerzelle}}
\end{aufgabe}

\uebersichtskasten{Zeit rechnet nicht in Zehnerschritten: 1 min = 60 s, 1 h = 60 min, 1 Tag = 24 h. \\ Dezimalstunde: Die Zahl hinter dem Komma sind Anteile einer Stunde, keine Minuten. 3,5 h = 3,5 · 60 min = 210 min. Nicht 3 h 50 min. \\ Uhrzeit plus Dauer: erst die vollen Stunden, dann die Minuten; ab 60 Minuten eine Stunde weiterzählen. 11:38 Uhr + 2 h = 13:38 Uhr, + 35 min = 14:13 Uhr.}
